% from ega4-14-fr.tex
\oldpage[IV-3]{221}

$z_i$ entiers sur $A$), que l'on peut supposer (en multipliant au besoin $s$ par un
élément de $A$) être des éléments du corps des fractions de $B_1$, entiers sur
$A_s$, de sorte que si $V$ est l'ouvert affine $D(s)=\operatorname{Spec}(A_s)\subset Y$, $Y'_1$
s'identifie à $Y_1\times_V V'$, où $V'$ est le spectre de la $A_s$-algèbre finie
engendrée par les $z_i/s$~; $g:V'\to V$ est donc un morphisme fini surjectif et
$g_1=g_{(Y_1)}$. De plus, appliquant la méthode de (8.1.2, $a)$), on peut supposer
que chacun des $U'_j$ est l'image réciproque par $p:Y'_1\to V'$ d'un ouvert $W'_j$
de $V'$, tel que les $W'_j$ recouvrent $V'$ (8.3.11), et que chacune des sections
$h'_j$ est de la forme $(v'_j)_{(Y_1)}$, où $v'_j$ est une $W'_j$-section de
$X\times_Y W'_j$ (8.8.2, (i)). On est donc bien ramené à prouver la proposition
lorsque $Y=\operatorname{Spec}(A)$, $A$ étant un anneau local intègre et noethérien. On sait alors
qu'il existe une $A$-algèbre intègre finie $B$, ayant même corps des fractions que
$A$, telle que $A\subset B$ et que $\operatorname{Spec}(B)$ soit géométriquement unibranche
(($\mathbf 0$, 23.2.5) et (6.15.5))~; comme le morphisme
$\operatorname{Spec}(B)\to\operatorname{Spec}(A)$ est surjectif, on peut évidemment remplacer $Y$ par
$\operatorname{Spec}(B)$ et $X$ par $X\otimes_A B$ pour prouver la proposition. Raisonnant comme
au début de la réduction III), on peut donc supposer $A$ local, intègre et
$Y=\operatorname{Spec}(A)$ géométriquement unibranche.

\emph{IV) Réduction au cas où $X$ est intègre, affine, et $f$ quasi-fini,
surjectif, birationnel et universellement ouvert.} -- Supposons donc $Y=\operatorname{Spec}(A)$
intègre, local et géométriquement unibranche. Il existe alors un sous-préschéma
irréductible $X_0$ de $X$ tel que la restriction $f_0:X_0\to Y$ de $f$ soit un
morphisme quasi-fini et dominant et que $f_0(X_0)$ contienne le point fermé $y$ de
$Y$ (14.5.9)~; comme $Y$ est géométriquement unibranche, il résulte de (14.4.1) que
$f_0$ est encore \emph{universellement ouvert}. Comme en outre on peut supposer
$X_0$ réduit, donc intègre, on voit qu'on peut, en remplaçant $X$ par $X_0$,
supposer que $X$ est \emph{intègre} et $f$ \emph{quasi-fini et dominant}, et tel que
$f(X)$ contienne $y$~; comme $f$ est \emph{ouvert} et que tout ouvert de $Y$
contenant $y$ est égal à $Y$, $f$ est surjectif.

Soient $\xi$, $\eta$ les points génériques de $X$ et $Y$ respectivement~;
$k(\xi)$ est donc une extension \emph{finie} de $K=k(\eta)$. Par suite (4.6.8), il
y a une extension \emph{finie} $K'$ de $K$ telle que
$(\operatorname{Spec}(k(\xi)\otimes_K K'))_{\mathrm{red}}$ soit géométriquement réduit sur $K$ et
que ses composantes irréductibles soient géométriquement irréductibles~; il en
résulte ((4.5.9) et (4.6.1)) que les corps résiduels de
$\operatorname{Spec}(k(\xi)\otimes_K K')$ sont des extensions finies, primaires et séparables de
$K'$, donc sont \emph{égaux à} $K'$. Appliquant de nouveau ($\mathbf 0$, 23.2.5) et
(6.15.5), il existe une $A$-algèbre intègre \emph{finie} $B$, ayant $K'$ pour
corps des fractions, contenant $A$ et telle que $\operatorname{Spec}(B)$ soit géométriquement
unibranche~; on peut donc de nouveau remplacer $Y$ par $\operatorname{Spec}(B)$ et $X$ par
$(X\otimes_A B)_{\mathrm{red}}$ pour prouver la proposition~; la réduction III)
permet alors de supposer encore $A$ local, intègre et $Y=\operatorname{Spec}(A)$ géométriquement
unibranche, de point fermé $y$. En outre, $f$ est alors un morphisme \emph{quasi-fini,
surjectif et universellement ouvert}~; chacune des composantes irréductibles $X_i$
de $X$ ($1\leq i\leq m$) domine $Y$ (1.10.4) et est \emph{birationnelle} sur $Y$.
Soit alors $x$ un point de $f^{-1}(y)$ et soit $X_i$ une composante irréductible de
$X$ contenant $x$~; désignons encore par $X_i$ le sous-préschéma réduit (donc
intègre) de $X$ ayant $X_i$ pour espace sous-jacent, et soit $f_i:X_i\to Y$ la
restriction de $f$. Appliquant de nouveau (14.4.1) au morphisme quasi-fini et
dominant $f_i$, on voit que $f_i$ est universellement ouvert~; si $U$ est un ouvert
affine de $X_i$ contenant $x$, $f_i(U)$ est donc

